\documentclass[conference]{IEEEtran}

% ── Paquetes ──────────────────────────────────────────────
\usepackage[utf8]{inputenc}
\usepackage[spanish]{babel}
\usepackage{graphicx}
\usepackage{float}
\usepackage{listings}
\usepackage{xcolor}
\usepackage{hyperref}
\usepackage{cite}
\usepackage{amsmath}
\usepackage{booktabs}

% ── Estilo de código terminal ─────────────────────────────
\lstdefinestyle{terminal}{
    backgroundcolor=\color{black},
    basicstyle=\ttfamily\scriptsize\color{white},
    breaklines=true,
    frame=single,
    rulecolor=\color{gray},
    showstringspaces=false,
    numbers=none,
    xleftmargin=3pt,
    xrightmargin=3pt,
}

\lstdefinestyle{cisco}{
    backgroundcolor=\color{black!90},
    basicstyle=\ttfamily\scriptsize\color{green!70!white},
    breaklines=true,
    frame=single,
    rulecolor=\color{gray},
    showstringspaces=false,
    numbers=none,
    xleftmargin=3pt,
    xrightmargin=3pt,
}

% ─────────────────────────────────────────────────────────
\begin{document}

% ── TÍTULO Y AUTORES ──────────────────────────────────────
\title{Lab No. 07 -- Basic Infrastructure and Network Layer}

\author{
    \IEEEauthorblockN{Sara Viviana Arteaga Rodríguez}
    \IEEEauthorblockA{Ingeniería de Sistemas}
    \and
    \IEEEauthorblockN{Diego Alejandro Mesa Alonso}
    \IEEEauthorblockA{Ingeniería de Sistemas}
    \and
    \IEEEauthorblockN{Julián Felipe Morales Zambrano}
    \IEEEauthorblockA{Ingeniería de Sistemas}
}

\maketitle

% ── RESUMEN ───────────────────────────────────────────────
\begin{abstract}
Este laboratorio cubre la configuración y administración de redes a nivel de infraestructura.
Se trabaja en tres sistemas operativos (Slackware, NetBSD y Windows Server) y se divide en
varias secciones principales. Primero se instalan y estudian herramientas de diagnóstico de red
como \texttt{netstat}, \texttt{vnstat}, \texttt{route} y \texttt{ethtool}, creando scripts con
menús interactivos que presentan la información de forma clara. También se instala software de
monitoreo de red basado en SNMP para supervisar las máquinas virtuales en tiempo real. En la
parte de experimentos, se analiza el comportamiento del protocolo ICMP usando traceroute hacia
diferentes sitios web alrededor del mundo, y se responden preguntas teóricas sobre comandos de
routers Cisco. Luego se realiza la configuración física de routers Cisco, incluyendo contraseñas,
nombre del dispositivo, sincronización de pantalla y mensaje del día. Se interconectan dos routers
mediante cable serial, se explican conceptos como DTE/DCE y \texttt{clock rate}, y finalmente se
configura enrutamiento estático para lograr comunicación completa entre todos los equipos del
curso, verificando con \texttt{ping} y \texttt{traceroute}.
\end{abstract}

\begin{IEEEkeywords}
redes, SNMP, Zabbix, routing estático, Cisco, ICMP, traceroute, NetBSD, Slackware, Windows Server
\end{IEEEkeywords}

% ─────────────────────────────────────────────────────────

% ── SECCIÓN I: OTHER USEFUL COMMANDS ─────────────────────
\section{Other Useful Commands}

\subsection{red\_info.sh -- NetBSD}

A continuación se muestra el menú principal del script de información de red ejecutado en NetBSD:

\begin{figure}[H]
    \centering
    \includegraphics[width=\columnwidth]{img_01_netbsd_menu}
    \caption{Menú principal de red\_info.sh en NetBSD.}
    \label{fig:netbsd_menu}
\end{figure}

Por ejemplo, al seleccionar la opción 3 (Estadísticas de interfaces):

\begin{figure}[H]
    \centering
    \includegraphics[width=\columnwidth]{img_02_netbsd_opcion3}
    \caption{Opción 3 -- Estadísticas de interfaces en NetBSD.}
    \label{fig:netbsd_op3}
\end{figure}

\subsection{Verificar\_puerto.sh -- NetBSD}

El script verificador de puertos indica si un puerto está abierto o cerrado e identifica
el servicio activo:

\begin{figure}[H]
    \centering
    \includegraphics[width=\columnwidth]{img_03_netbsd_puerto_abierto}
    \caption{Puerto 80 abierto -- Verificar\_puerto.sh en NetBSD.}
    \label{fig:netbsd_p80}
\end{figure}

\begin{figure}[H]
    \centering
    \includegraphics[width=\columnwidth]{img_04_netbsd_puerto_cerrado}
    \caption{Puerto 8080 cerrado -- Verificar\_puerto.sh en NetBSD.}
    \label{fig:netbsd_p8080}
\end{figure}

\subsection{red\_info.sh -- Slackware}

\begin{figure}[H]
    \centering
    \includegraphics[width=\columnwidth]{img_05_slackware_menu}
    \caption{Menú principal de red\_info.sh en Slackware.}
    \label{fig:slack_menu}
\end{figure}

\subsection{Verificar\_puerto.sh -- Slackware}

\begin{figure}[H]
    \centering
    \includegraphics[width=\columnwidth]{img_06_slackware_puerto}
    \caption{Puerto 22 abierto (SSH) -- Verificar\_puerto.sh en Slackware.}
    \label{fig:slack_puerto}
\end{figure}

\subsection{red\_info.sh -- Windows Server}

\begin{figure}[H]
    \centering
    \includegraphics[width=\columnwidth]{img_07_windows_red_info_notepad}
    \caption{Apertura de red\_info.ps1 en Windows Server.}
    \label{fig:win_notepad}
\end{figure}

\begin{figure}[H]
    \centering
    \includegraphics[width=\columnwidth]{img_08_windows_menu}
    \caption{Menú principal de red\_info.ps1 en Windows Server.}
    \label{fig:win_menu}
\end{figure}

\subsection{Verificar\_puerto.sh -- Windows Server}

\begin{figure}[H]
    \centering
    \includegraphics[width=\columnwidth]{img_09_windows_puerto}
    \caption{Puerto 80 abierto -- Verificar\_puerto.ps1 en Windows Server.}
    \label{fig:win_puerto}
\end{figure}

% ── SECCIÓN II: NETWORK MANAGEMENT ───────────────────────
\section{Network Management}

Prueba visual de que la red está siendo monitoreada en tiempo real desde un solo punto
central con Zabbix.

\begin{figure}[H]
    \centering
    \includegraphics[width=\columnwidth]{img_10_zabbix_hosts}
    \caption{Panel de hosts en Zabbix.}
    \label{fig:zabbix_hosts}
\end{figure}

A continuación se muestra cómo Zabbix monitorea el disco, CPU, memoria y red de cada
sistema operativo en tiempo real:

\begin{figure}[H]
    \centering
    \includegraphics[width=\columnwidth]{img_11_zabbix_netbsd}
    \caption{Monitoreo de NetBSD en tiempo real con Zabbix.}
    \label{fig:zabbix_netbsd}
\end{figure}

\begin{figure}[H]
    \centering
    \includegraphics[width=\columnwidth]{img_12_zabbix_slackware}
    \caption{Monitoreo de Slackware en tiempo real con Zabbix.}
    \label{fig:zabbix_slack}
\end{figure}

\begin{figure}[H]
    \centering
    \includegraphics[width=\columnwidth]{img_13_zabbix_windows}
    \caption{Monitoreo de Windows Server en tiempo real con Zabbix.}
    \label{fig:zabbix_win}
\end{figure}

Se ejecutó \texttt{snmpconf -g basic\_setup}, un asistente interactivo para configurar
el demonio SNMP. Se configuraron 3 secciones:

\begin{enumerate}
    \item \textbf{Información del sistema:} Se definió la ubicación física
    \textit{``Laboratorio de Redes''} y el contacto \textit{``sara@universidad.edu''}.

    \item \textbf{Control de acceso (admin):} Usuario \texttt{admin} con nivel
    \texttt{auth}, con permisos de lectura y escritura.

    \item \textbf{Control de acceso (monitor):} Usuario \texttt{monitor} con nivel
    \texttt{auth}, solo lectura. Es el usuario que Zabbix usa para consultar el sistema.
\end{enumerate}

SNMP permite a Zabbix consultar remotamente uso de CPU, memoria, disco, interfaces de
red y procesos activos. Sin esta configuración, Zabbix no podría obtener datos del
servidor vía SNMP.

\begin{figure}[H]
    \centering
    \includegraphics[width=\columnwidth]{img_14_snmpconf}
    \caption{Configuración del demonio SNMP con snmpconf -g basic\_setup.}
    \label{fig:snmpconf}
\end{figure}

Con \texttt{zabbix\_server onestatus}, \texttt{zabbix\_agentd onestatus} y
\texttt{snmpd onestatus} se confirmó que todos los servicios están activos y
respondiendo correctamente en el puerto 161.

\begin{figure}[H]
    \centering
    \includegraphics[width=\columnwidth]{img_15_zabbix_status}
    \caption{Estado de zabbix\_server, zabbix\_agentd y snmpd.}
    \label{fig:zabbix_status}
\end{figure}

% ── SECCIÓN III: NETWORK ADMINISTRATION - AZURE ──────────
\section{Network Administration -- Azure}

Se accedió al portal \url{https://portal.azure.com} con la cuenta institucional, se
navegó a \textit{Education} $\rightarrow$ \textit{Templates} y se desplegó la plantilla
\textbf{Web App Deployment from GitHub}.

\begin{figure}[H]
    \centering
    \includegraphics[width=\columnwidth]{img_17_azure_deploy}
    \label{fig:azure_deploy}
\end{figure}

\begin{figure}[H]
    \centering
    \includegraphics[width=\columnwidth]{17.2.png}
    \label{fig:azure_deploy}
\end{figure}

\begin{figure}[H]
    \centering
    \includegraphics[width=\columnwidth]{17.3png.png}
    \label{fig:azure_deploy}
\end{figure}

\textbf{¿Qué pasa al refrescar?} Se incrementa el número de solicitudes en tiempo real,
reflejado en los gráficos de Live Metrics. También varían el tiempo de respuesta y el uso
de recursos.

\textbf{¿Qué muestra?} Solicitudes, rendimiento, tiempos de respuesta, errores y
actividad de usuarios en la aplicación.

\textbf{¿Para qué sirve?} Permite monitorear el estado de la aplicación, detectar fallos
y optimizar el rendimiento en tiempo real.

\textbf{Capa de red:} Permite la comunicación entre cliente y servidor mediante
direcciones IP y rutas.

\textbf{Transporte y aplicación:} La capa de transporte asegura la entrega correcta
(TCP), mientras que la capa de aplicación interpreta las solicitudes (HTTP).

% ── SECCIÓN IV: USE OF ICMP MESSAGES ─────────────────────
\section{Use of ICMP Messages}

\subsection{Traceroute hacia Stanford y Computer Science Laboratory}

\begin{figure}[H]
    \centering
    \includegraphics[width=\columnwidth]{img_18_traceroute_stanford}
    \caption{Resultado de traceroute hacia Stanford University.}
    \label{fig:stanford}
\end{figure}

\begin{figure}[H]
    \centering
    \includegraphics[width=\columnwidth]{img_19_traceroute_cs}
    \caption{Resultado de traceroute hacia el Computer Science Laboratory.}
    \label{fig:cs}
\end{figure}

\subsection{Traceroute hacia una página en Francia}

\begin{figure}[H]
    \centering
    \includegraphics[width=\columnwidth]{img_20_traceroute_francia}
    \caption{Traceroute hacia página en Francia.}
    \label{fig:francia}
\end{figure}

\subsection{Open Visual Traceroute}

\begin{figure}[H]
    \centering
    \includegraphics[width=\columnwidth]{img_22_ovt_prueba}
    \caption{Prueba inicial de Open Visual Traceroute.}
    \label{fig:ovt_prueba}
\end{figure}

\subsection{Fabricantes de automóviles en diferentes continentes}

Se realizaron trazas hacia sitios web de fabricantes en distintos continentes:
\texttt{toyota.jp} (Japón), \texttt{bmw.de} (Alemania), \texttt{ford.com}
(Estados Unidos), \texttt{hyundai.com/au} (Australia) y \texttt{renault.fr} (Francia).

\begin{figure}[H]
    \centering
    \includegraphics[width=\columnwidth]{img_23_toyota}
    \caption{Open Visual Traceroute -- toyota.jp (Japón).}
    \label{fig:toyota}
\end{figure}

\begin{figure}[H]
    \centering
    \includegraphics[width=\columnwidth]{img_24_bmw}
    \caption{Open Visual Traceroute -- bmw.de (Alemania).}
    \label{fig:bmw}
\end{figure}

\begin{figure}[H]
    \centering
    \includegraphics[width=\columnwidth]{img_25_ford}
    \caption{Open Visual Traceroute -- ford.com (Estados Unidos).}
    \label{fig:ford}
\end{figure}

\begin{figure}[H]
    \centering
    \includegraphics[width=\columnwidth]{img_26_hyundai}
    \caption{Open Visual Traceroute -- hyundai.com/au (Australia).}
    \label{fig:hyundai}
\end{figure}

\begin{figure}[H]
    \centering
    \includegraphics[width=\columnwidth]{img_27_renault}
    \caption{Open Visual Traceroute -- renault.fr (Francia).}
    \label{fig:renault}
\end{figure}

% ── SECCIÓN V: ROUTER COMMANDS ───────────────────────────
\section{Some Questions About Router Commands}

\begin{enumerate}
    \item \textbf{enable password vs enable secret:}
    \texttt{Enable password} guarda la contraseña en texto plano y \texttt{enable secret}
    la cifra con MD5. Si ambas están configuradas, \texttt{enable secret} tiene precedencia
    y \texttt{enable password} es ignorado.

    \item \textbf{Console vs VTY:}
    \texttt{Console} es el acceso físico directo al router mediante cable de consola
    (RJ-45/DB9). \texttt{VTY} son las líneas virtuales para acceso remoto por Telnet o SSH,
    normalmente de 0 a 4 (5 sesiones simultáneas).

    \item \textbf{Proceso de arranque:}
    El router ejecuta el POST, carga el Bootstrap desde ROM, busca y carga el IOS desde
    Flash, luego carga la \texttt{startup-config} desde NVRAM a la RAM como
    \texttt{running-config}. Si no existe, entra al Setup Dialog.

    \item \textbf{Tipos de memoria:}
    ROM, Flash Memory, NVRAM y DRAM.

    \item \textbf{startup-config vs running-config:}
    \texttt{startup-config} está en NVRAM y persiste al apagar el router.
    \texttt{running-config} está en RAM y es la configuración activa; se pierde si el
    router se apaga sin guardar.
\end{enumerate}

% ── SECCIÓN VI: BASIC CONFIGURATION OF ROUTERS ───────────
\section{Setup: Access and Basic Configuration of Routers}

Se realizó la configuración básica del router incluyendo contraseñas, nombre, sincronización
de pantalla, descripción de interfaces, desactivación del DNS lookup y mensaje del día:

\begin{figure}[H]
    \centering
    \includegraphics[width=\columnwidth]{enable.png}
    \label{fig:azure_deploy}
\end{figure}
\begin{figure}[H]
    \centering
    \includegraphics[width=\columnwidth]{configT.png}
    \label{fig:azure_deploy}
\end{figure}
\begin{figure}[H]
    \centering
    \includegraphics[width=\columnwidth]{saraconfig.png}
    \label{fig:azure_deploy}
\end{figure}
Se configuraron las interfaces y se verificó conectividad entre computadores:
\begin{figure}[H]
    \centering
    \includegraphics[width=\columnwidth]{connectivity.png}
    \label{fig:azure_deploy}
\end{figure}

\begin{figure}[H]
    \centering
    \includegraphics[width=\columnwidth]{ping1.png}
    \label{fig:azure_deploy}
\end{figure}
\begin{figure}[H]
    \centering
    \includegraphics[width=\columnwidth]{ping2.png}
    \label{fig:azure_deploy}
\end{figure}
% ── SECCIÓN VII: SERIAL INTERCONNECTION ──────────────────
\section{Setup -- Serial Interconnection}

\textbf{¿Qué es un null modem?}
Es un cable serial que conecta dos dispositivos directamente sin necesidad de un módem real,
cruzando las líneas TX y RX para que dos equipos puedan comunicarse punto a punto. En
laboratorio se usa para simular una conexión WAN serial entre dos routers.

\textbf{¿Para qué sirve \texttt{clock rate}?}
En conexiones seriales, el dispositivo DCE debe generar la señal de reloj que sincroniza la
transmisión. El comando \texttt{clock rate} establece esa velocidad (ej. \texttt{clock rate 64000});
sin ella la interfaz serial no puede intercambiar tráfico.

\textbf{DTE y DCE:}
DTE (Data Terminal Equipment) recibe la señal de reloj; DCE (Data Communications Equipment)
la genera. En el laboratorio, el extremo DCE del cable serial debe configurar el
\texttt{clock rate}. Para identificarlo se usa:

\begin{lstlisting}[style=cisco]
Router# show controllers serial 0/0/0
\end{lstlisting}

% ── SECCIÓN VIII: STATIC ROUTING ─────────────────────────
\section{Static Routing}
\begin{figure}[H]
    \centering
    \includegraphics[width=\columnwidth]{img_29_static_routing}
    \caption{Configuración de rutas estáticas en los routers.}
    \label{fig:static}
\end{figure}

% ── SECCIÓN IX: CLOSURE ───────────────────────────────────
\section{Closure}

Se borró la configuración de los routers con \texttt{erase startup-config} seguido de
\texttt{reload}, dejando los dispositivos sin contraseñas ni rutas configuradas y listos
para una nueva sesión. Se organizaron cables y equipos en el laboratorio.

% ── CONCLUSIÓN ────────────────────────────────────────────
\section{Conclusión}

Este laboratorio permitió aplicar de manera práctica conceptos fundamentales de redes que
normalmente se estudian de forma teórica. Trabajar con tres sistemas operativos distintos
---Slackware, NetBSD y Windows Server--- para desarrollar scripts de diagnóstico de red
evidenció cómo una misma tarea puede requerir enfoques completamente diferentes según la
plataforma, y reforzó la capacidad de interpretar y presentar información del sistema de
forma clara y automatizada.

Configurar un servidor de monitoreo con SNMP fue especialmente útil para entender por qué
la observabilidad es un componente crítico en redes empresariales. Poder ver en tiempo real
el estado de CPU, memoria y tráfico de las máquinas virtuales deja claro que sin estas
herramientas, detectar y resolver problemas en una red sería mucho más difícil. En esa misma
línea, explorar Application Insights en Azure mostró cómo esta necesidad de visibilidad también
existe en la nube, donde las aplicaciones deben monitorearse constantemente para garantizar
su rendimiento.

La parte más práctica del laboratorio fue sin duda la configuración de routers Cisco.
Recuperar acceso, establecer contraseñas, conectar equipos por puerto serial y definir rutas
estáticas manualmente ayudó a comprender de forma concreta cómo se construye y administra
una red real. Verificar la conectividad con ping y traceroute al final del proceso confirmó
que cada decisión de configuración tiene un impacto directo en la comunicación entre los
dispositivos.

% ── ANEXOS ────────────────────────────────────────────────
\section*{Anexos}

\subsection*{Instalación de Zabbix en NetBSD}
\begin{figure}[H]
    \centering
    \includegraphics[width=\columnwidth]{Net1.png}
    \label{fig:azure_deploy}
\end{figure}
\begin{figure}[H]
    \centering
    \includegraphics[width=\columnwidth]{Net2.png}
    \label{fig:azure_deploy}
\end{figure}
\begin{figure}[H]
    \centering
    \includegraphics[width=\columnwidth]{Net3.png}
    \label{fig:azure_deploy}
\end{figure}
\begin{figure}[H]
    \centering
    \includegraphics[width=\columnwidth]{Net4.png}
    \label{fig:azure_deploy}
\end{figure}
\begin{figure}[H]
    \centering
    \includegraphics[width=\columnwidth]{Net5.png}
    \label{fig:azure_deploy}
\end{figure}
\begin{figure}[H]
    \centering
    \includegraphics[width=\columnwidth]{Net6.png}
    \label{fig:azure_deploy}
\end{figure}
\begin{figure}[H]
    \centering
    \includegraphics[width=\columnwidth]{Net7.png}
    \label{fig:azure_deploy}
\end{figure}
\begin{figure}[H]
    \centering
    \includegraphics[width=\columnwidth]{Net8.png}
    \label{fig:azure_deploy}
\end{figure}
\begin{figure}[H]
    \centering
    \includegraphics[width=\columnwidth]{Net9.png}
    \label{fig:azure_deploy}
\end{figure}
\begin{figure}[H]
    \centering
    \includegraphics[width=\columnwidth]{Net10.png}
    \label{fig:azure_deploy}
\end{figure}
\begin{figure}[H]
    \centering
    \includegraphics[width=\columnwidth]{Net11.png}
    \label{fig:azure_deploy}
\end{figure}
\begin{figure}[H]
    \centering
    \includegraphics[width=\columnwidth]{Net12.png}
    \label{fig:azure_deploy}
\end{figure}
\begin{figure}[H]
    \centering
    \includegraphics[width=\columnwidth]{Net13.png}
    \label{fig:azure_deploy}
\end{figure}
\begin{figure}[H]
    \centering
    \includegraphics[width=\columnwidth]{Net14.png}
    \label{fig:azure_deploy}
\end{figure}
\begin{figure}[H]
    \centering
    \includegraphics[width=\columnwidth]{Net15.png}
    \label{fig:azure_deploy}
\end{figure}
\begin{figure}[H]
    \centering
    \includegraphics[width=\columnwidth]{Net16.png}
    \label{fig:azure_deploy}
\end{figure}
\begin{figure}[H]
    \centering
    \includegraphics[width=\columnwidth]{Net17.png}
    \label{fig:azure_deploy}
\end{figure}
\begin{figure}[H]
    \centering
    \includegraphics[width=\columnwidth]{Net18.png}
    \label{fig:azure_deploy}
\end{figure}

\subsection*{Instalación de Zabbix en Slackware}

\begin{figure}[H]
    \centering
    \includegraphics[width=\columnwidth]{Slack1.png}
    \label{fig:azure_deploy}
\end{figure}
\begin{figure}[H]
    \centering
    \includegraphics[width=\columnwidth]{Slack2.png}
    \label{fig:azure_deploy}
\end{figure}
\begin{figure}[H]
    \centering
    \includegraphics[width=\columnwidth]{Slack3.png}
    \label{fig:azure_deploy}
\end{figure}
\begin{figure}[H]
    \centering
    \includegraphics[width=\columnwidth]{Slack4.png}
    \label{fig:azure_deploy}
\end{figure}
\subsection*{Instalación de Zabbix en Windows Server}

\begin{figure}[H]
    \centering
    \includegraphics[width=\columnwidth]{Win1.png}
    \label{fig:azure_deploy}
\end{figure}
\begin{figure}[H]
    \centering
    \includegraphics[width=\columnwidth]{Win2.png}
    \label{fig:azure_deploy}
\end{figure}

\end{document}